\section{Numerical Experiments}
\label{sec:experiments}

This section evaluates the WAIT and Nested WAIT algorithms across a range of workloads, comparing them against two widely deployed inference schedulers: vLLM~\citep{kwon2023efficient}, which prioritizes new arrivals under a first-come-first-serve (FCFS) policy, and Sarathi~\citep{agrawal2023sarathi}, which combines chunked prefill with a strict admission cap. Our presentation is organized around the central diagnostic suggested by queueing theory: how mean end-to-end latency scales with the arrival rate, and where each policy's stability region ends. Section~\ref{sec:exp_setup} describes the simulator, hardware, workloads, and parameters. Section~\ref{sec:exp_stability} presents the main mean-latency-versus-rate comparison and reports the resulting stability orderings.
%% TODO when later sections are written:
%% Section~\ref{sec:exp_known} provides additional rate-sweep results for WAIT in the known-output setting.
%% Section~\ref{sec:exp_unknown} extends to Nested WAIT with unknown output lengths and to the lmsys real workload.
%% Section~\ref{sec:exp_eviction} isolates the near-boundary regime and the eviction-cascade behavior.

\subsection{Experimental Setup}
\label{sec:exp_setup}

\paragraph{Simulator and hardware.}
We run all experiments on Vidur~\citep{agrawal2024vidur}, a discrete-event simulator that has been calibrated against profiling traces collected on real GPUs. The simulated platform is a single NVIDIA A100~80\,GB serving Llama-2-7B with tensor- and pipeline-parallelism degree one and a KV-cache block size of 16~tokens. To keep the simulator's iteration-time predictions faithful to real hardware over the regime used in our experiments, we additionally validate Vidur against direct vLLM measurements on a physical A100; the comparison is shown as Figure~\ref{fig:sim_fidelity} in Appendix~\ref{app:sim_fidelity}, and we have replicated the WAIT mechanism on a production SGLang deployment as a complementary sanity check (Appendix~\ref{app:sglang}).

\paragraph{Algorithms compared.}
We compare three policies: vLLM, Sarathi, and the WAIT (or Nested WAIT) algorithm of Algorithm~\ref{alg:wait} (resp.\ Algorithm~\ref{alg:nested_wait}). Throughout, we refer to the single-type setting (with known output length) as the WAIT setting, and to the multi-type and unknown-output settings as the Nested WAIT setting. Sarathi is configured with chunk size $512$ and an admission cap of $512$; vLLM uses a watermark of $0.01$ on the KV-cache pool. WAIT uses chunk size $256$ and threshold $\sum_k n_k = 21$; Nested WAIT uses configuration $(\mathrm{cs}, \sum_k n_k)$ that varies with workload, as specified in Table~\ref{tab:exp_setup}. All three policies share the same simulator, hardware, and KV-cache settings.

\paragraph{Workloads.}
We consider three classes of workloads, each with Poisson arrivals at rate $\lambda$:
\begin{itemize}
\item \textbf{Single-type (Section~\ref{sec:exp_stability})}: prefill length $512$ and decode length $20$ (denoted \texttt{p512d20}); this is the canonical workload for the WAIT setting and isolates the effect of the threshold mechanism.
\item \textbf{Multi-type W3 (Section~\ref{sec:exp_stability})}: a $70/30$ mixture of \texttt{p512d20} and \texttt{p512d50}; the two types share the prefill profile but differ in decode length, which exercises the segment design of Nested WAIT.
\item \textbf{Real workload}: $50{,}000$ prompts sampled from the lmsys-chat-1m dataset~\citep{zheng2023lmsys}; Section~\ref{sec:exp_unknown} reports results on this workload.
\end{itemize}

\paragraph{Memory safeguard.}
The threshold $\sum_k n_k$ used by WAIT and Nested WAIT is an explicit upper bound on the number of in-system requests, which in turn caps the KV-cache footprint at $(\sum_k n_k) \cdot M^*_{\max}$ where $M^*_{\max}$ is the largest per-request memory profile across active types. Because this bound is enforced by construction at every batching decision, no eviction is required to satisfy the memory constraint $C \geq M^*$ assumed in Sections~\ref{sec:wait} and~\ref{sec:nested_wait}. By contrast, vLLM and Sarathi rely on watermarks and on-demand eviction to keep within the same constraint, which has consequences for stability that we examine next.

\paragraph{Parameters.}
Table~\ref{tab:exp_setup} summarizes the parameters used in each workload, including the batch budget $B$, the equilibrium memory $M^*$, the capacity $C$ (in KV-cache blocks), the iteration-time intercept and slope $(d_0, d_1)$ extracted from Vidur's profile of A100 + Llama-2-7B, and the threshold $\sum_k n_k$ used by WAIT and Nested WAIT.
%% TODO: fill this table once §6.3-6.6 are finalized; the entries are
%% available in CLAUDE.md but kept in one consolidated table to
%% address Reviewer 2's request for a single parameter summary.
\begin{table}[ht]
\centering
\begin{tabular}{lccccc}
\toprule
Workload & $B$ & $M^*$ (blocks) & $C$ (blocks) & $(d_0, d_1)$ & $\sum_k n_k$ \\
\midrule
Single-type \texttt{p512d20} & 512 & TBD & 29952 & TBD & 21 \\
Multi-type W3 & 512 & TBD & 29952 & TBD & 20--30 (rate-dependent) \\
Real (lmsys) & 512 & TBD & 29952 & TBD & TBD \\
\bottomrule
\end{tabular}
\caption{Parameters used in the experiments.}
\label{tab:exp_setup}
\end{table}

\subsection{Stability Region: Mean Latency versus Arrival Rate}
\label{sec:exp_stability}

In an open queueing system, throughput equals the arrival rate $\lambda$ in any rate region where the policy keeps the queue positive recurrent (i.e.\ stable); throughput differences between policies therefore reflect the size of their respective stability regions. The standard diagnostic is to plot mean end-to-end latency as a function of $\lambda$ and read the divergence point off each curve. We adopt this diagnostic throughout this section.

\paragraph{Single-type (WAIT).}
Figure~\ref{fig:single_type_rate} reports the result for the canonical single-type workload. The left panel shows the mean end-to-end latency on a logarithmic scale; the right panel shows the corresponding effective throughput, defined as the completion rate that each policy is able to sustain. As $\lambda$ increases, every policy eventually leaves its stability region, and the mean latency curve diverges. The WAIT algorithm sustains bounded latency up to $\lambda \approx 23$, whereas Sarathi loses stability near $\lambda \approx 22$ and vLLM near $\lambda \approx 19$. At every rate point in the swept range $\lambda \in \{12, 13, \dots, 24\}$, WAIT attains the smallest mean latency among the three policies; the gap relative to Sarathi widens monotonically from approximately $2\%$ at $\lambda = 12$ to over $70\%$ at $\lambda = 23$. The right panel reports the same ordering in throughput terms: WAIT remains on the ideal $\lambda$-line further than the two baselines and saturates at the highest sustainable throughput. The fact that the divergence points are well-separated, even though the three policies share the same hardware and memory budget, demonstrates that the threshold-based admission rule of WAIT enlarges the stability region beyond what the admission heuristics of vLLM and Sarathi achieve.

\begin{figure}[htbp]
    \centering
    \includegraphics[width=\textwidth]{Experiments_pdf/figure_B_single_type.pdf}
    \caption{Single-type workload (prefill $=512$, decode $=20$, Poisson arrivals): mean end-to-end latency (left, log scale) and effective throughput (right) as a function of arrival rate $\lambda$. \textbf{Left:} as $\lambda$ increases, each policy eventually enters an unstable regime where mean latency diverges; the divergence point delineates its stability region. WAIT sustains bounded latency up to $\lambda \approx 23$, whereas Sarathi and vLLM lose stability near $\lambda \approx 22$ and $\lambda \approx 19$, respectively. \textbf{Right:} the same ordering in terms of effective throughput (estimated from the last $\lambda$ at which each policy remains stable), with the dotted line indicating the ideal $\lambda = $~completion rate.}
    \label{fig:single_type_rate}
\end{figure}

\paragraph{Multi-type (Nested WAIT).}
%% TODO: insert Figure C and a parallel paragraph once the W3 vLLM
%% backfill (rate sweep 12-36, nreq=20000) finishes.
\textit{[Figure~C and accompanying paragraph for the Multi-type W3 workload will be inserted once the vLLM backfill at $\lambda \in \{12, \dots, 36\}$ completes.]}

\paragraph{Real workload (Nested WAIT).}
%% TODO: insert Figure D and accompanying paragraph after Step 4 (lmsys
%% rate sweep) is rendered; rate axis runs 10-150 QPS.
\textit{[Figure~D and accompanying paragraph for the lmsys-chat-1m workload will follow.]}

\paragraph{Finite-horizon scaling check.}
%% TODO: Figure E (nreq scaling 2k -> 20k) and short paragraph
%% confirming that the divergence points reported above are not
%% finite-horizon artifacts.

%% =====================================================================
%% Sections 6.3 -- 6.6 placeholders (to be drafted as figures complete)
%% =====================================================================

%% \subsection{Known Output: Detailed WAIT Results}
%% \label{sec:exp_known}
%% (Figures F-a, F-b, G-a, G-b)

%% \subsection{Unknown Output: Detailed Nested WAIT Results}
%% \label{sec:exp_unknown}
%% (Figures I-a, I-b, J-a, J-b; lmsys + long-decode)

%% \subsection{Near-Boundary Behavior and Eviction Cascade}
%% \label{sec:exp_eviction}
%% (Figures K, L; memory-constrained margin=0.6 r=4.0)
